\scanpage{49}
\mspage{\handwritten{\unsure{$\sim$|2}}}

% printed form
\deleted{\texttt{/*JOB=CRAY2,ACCT=*********,ALL=M06,MSGL=N}\\
\texttt{// JOB MSGLEVEL=(1,1),MSGCLASS=A}\\
\texttt{// EXEC PGM=IEFBR14}}

\handwritten{\uline{A examiner} : a) Extension de la conjecture fondamentale en géom. alg. anabélienne aux courbes anabéliennes \deleted{\unsure{stabl|stab|stables}} singulières (stables, avec propagateur normal). \unsure{?|P}}

\handwritten{b) Pour une $Y$, anabélienne singulière donnée (sur corps $\deleted{\unsure{k}}K$ de t.f. sur $\mathbb Q$), \deleted{\unsure{On}}est-il vrai qu'on récupère les \unsure{él.|cl.} de \deleted{(pts}$\mathrm{Sing}(Y)$ $\cup$ $\underbrace{\text{Pts marqués}}_{Y_\infty}$ en termes des \inserted{classes de conj.} sous-groupes de $\pi_1(Y)$ isom. à $\widehat{\mathbb Z}$, et \unsure{esst|exst|ext.} normalisés par $\pi_1(K)$ (l'opération de $\pi_1(K)$ dessus se fera automatiquement via le caractère cyclotomique)}

\handwritten{c) S'inspirant des compactifications MD des $M_{g,\nu}$, définir ce que c'est qu'une compactification anabélienne d'une variété lisse anabélienne (sur un corps $K$ de car. 0, qu'on peut supposer être de type fini sur $\mathbb Q$). Comportement par spécialisation ???}

% archive sheet number: 51
